\documentclass{article}

\usepackage{amsmath,amssymb}
\usepackage{xurl}
\usepackage[hidelinks]{hyperref}
\setlength{\emergencystretch}{2em}

\input{command}

\title{The Structure Theorem for Finite-Dimensional $C^\ast$-Algebras:
    Matrix-Realized Single-Block Scope}
\author{}
\date{2026-09-03}

\begin{document}
\maketitle
\gapnote{scope-restriction}{resolved}

\section{The source statement}

Schumacher and Werner state the structure theorem for finite-dimensional
$C^\ast$-algebras as their Proposition~``Csform''
\cite[lines~2082--2098]{SchumacherWerner2004}: every finite-dimensional
$C^\ast$-algebra $\mathcal A$ is determined up to isomorphism by a finite
sequence $n_1\geq n_2\geq\cdots\geq 1$ with
\begin{align}
    \mathcal A
        &\cong\bigoplus_\mu\MN{n_\mu}.
    \label{eq:sw04_csform_matrix_scope_source_decomposition}
\end{align}
The proof sketch given there cuts $\mathcal A$ by the minimal projections
$z_\mu$ of its centre,
\begin{align}
    \mathcal A
        &=\bigoplus_\mu z_\mu\mathcal A,
    \label{eq:sw04_csform_matrix_scope_central_cut}
\end{align}
and asserts that each summand $z_\mu\mathcal A$, whose centre consists of the
multiples of its unit, is isomorphic to a full matrix algebra $\MN{n_\mu}$.

Gross, Nesme, Vogts and Werner use exactly this building-block assertion, not
(\ref{eq:sw04_csform_matrix_scope_source_decomposition}) itself: a support
algebra $\mathcal R_x$ with trivial centre ``must be isomorphic to
$M_{r(x)}$''\cite[lines~1278--1282]{GrossNesmeVogtsWerner2012}.

\section{The point at issue}

Two features of
(\ref{eq:sw04_csform_matrix_scope_source_decomposition}) are outside the
formalized statement.

First, $\mathcal A$ is an \emph{abstract} finite-dimensional $C^\ast$-algebra.
The formalized statement quantifies over a $*$-subalgebra
\begin{align}
    S
        &\subseteq\MN{D},
    \label{eq:sw04_csform_matrix_scope_concrete_subalgebra}
\end{align}
that is, over an algebra already presented with a faithful finite-dimensional
$*$-representation. Every finite-dimensional $C^\ast$-algebra admits such a
presentation, but the passage from the abstract algebra to
(\ref{eq:sw04_csform_matrix_scope_concrete_subalgebra}) is the
Gelfand--Naimark theorem, which is not formalized here.

Second, the formalized statement is the single-block case. The centre of $S$ is
assumed to be the scalar image,
\begin{align}
    Z(S)
        &=\C\cdot\one,
    \label{eq:sw04_csform_matrix_scope_scalar_centre}
\end{align}
which is the hypothesis satisfied by one summand $z_\mu\mathcal A$ of
(\ref{eq:sw04_csform_matrix_scope_central_cut}). Under
(\ref{eq:sw04_csform_matrix_scope_scalar_centre}) the conclusion is a single
full matrix algebra,
\begin{align}
    S
        &\cong\MN{r},
        \qquad r\geq 1,
    \label{eq:sw04_csform_matrix_scope_single_block}
\end{align}
as a $*$-isomorphism.\footnote{Formal declaration:
\leanid{StarSubalgebra.exists_starAlgEquiv_matrix_of_isCentral} in
\doclink{QICLean/Algebra/StarSubalgebraFactor}.}
The decomposition
(\ref{eq:sw04_csform_matrix_scope_central_cut}) into minimal central
projections, and with it the passage from
(\ref{eq:sw04_csform_matrix_scope_single_block}) to
(\ref{eq:sw04_csform_matrix_scope_source_decomposition}), is not formalized.

\section{The formalized route}

The formal proof does not follow the minimal-projection sketch. It uses the
Wedderburn--Artin decomposition of $S$ as an abstract semisimple ring over the
algebraically closed field $\C$,
\begin{align}
    S
        &\cong\prod_{k=1}^{m}\MN{d_k},
        \qquad d_k\geq 1,
    \label{eq:sw04_csform_matrix_scope_wedderburn_product}
\end{align}
which is available because a $*$-subalgebra of $\MN{D}$ has vanishing Jacobson
radical.\footnote{Formal declarations: \leanid{StarSubalgebra.isSemisimpleRing}
and \leanid{StarSubalgebra.exists\_algEquiv\_pi\_matrix} in
\doclink{QICLean/Algebra/StarSubalgebraSemisimple}.}
Under (\ref{eq:sw04_csform_matrix_scope_scalar_centre}) the coordinate
idempotent of the $k$-th factor of
(\ref{eq:sw04_csform_matrix_scope_wedderburn_product}) is central, hence a
scalar multiple of the unit; a scalar that is the unit in one coordinate and
zero in another forces $m\leq 1$, and $m=0$ is excluded because $S$ contains
the unit of $\MN{D}$. So $m=1$ and $S$ is a simple ring, and the Wedderburn
theorem for simple algebras over $\C$ gives an algebra isomorphism
\begin{align}
    \varphi
        &:S\longrightarrow\MN{r}.
    \label{eq:sw04_csform_matrix_scope_algebra_isomorphism}
\end{align}

The isomorphism (\ref{eq:sw04_csform_matrix_scope_algebra_isomorphism}) is
supplied by ring theory alone and carries no compatibility with the adjoint.
The formal development therefore adds the step that the source leaves inside
its phrase ``are then seen to be isomorphic'': transporting the adjoint of $S$
through $\varphi$ gives an algebra automorphism of $\MN{r}$, which by
Skolem--Noether is conjugation by an invertible $H$, so that
\begin{align}
    \varphi(y^\dagger)
        &=H^{-1}\varphi(y)^\dagger H,
        \qquad y\in S.
    \label{eq:sw04_csform_matrix_scope_transported_adjoint}
\end{align}
Applying (\ref{eq:sw04_csform_matrix_scope_transported_adjoint}) twice shows
that $H^\dagger$ is a scalar multiple of $H$. Evaluating the trace of
$\varphi^{-1}$ on the rank-one matrices $Y=vw^\dagger$, for which
\begin{align}
    Y^\dagger HY
        &=(v^\dagger Hv)\,ww^\dagger,
    \label{eq:sw04_csform_matrix_scope_rank_one_sandwich}
\end{align}
shows that $v^\dagger Hv$ never vanishes for $v\neq 0$ and that, after $H$ is
rescaled by one scalar, the Hermitian form has constant sign. The rescaled
matrix is positive definite; writing it as $R^\dagger R$ with $R$ invertible,
the corrected map
\begin{align}
    y
        &\longmapsto R\,\varphi(y)\,R^{-1}
    \label{eq:sw04_csform_matrix_scope_corrected_isomorphism}
\end{align}
preserves adjoints and realizes
(\ref{eq:sw04_csform_matrix_scope_single_block}).\footnote{Formal declaration:
\leanid{StarSubalgebra.exists\_starAlgEquiv\_of\_algEquiv\_matrix}.}

\section{Verdict}

\emph{Closed for the matrix-realized single-block case.} The formal statement
proves (\ref{eq:sw04_csform_matrix_scope_single_block}) for a $*$-subalgebra of
a nonzero full matrix algebra whose centre is scalar, which is the assertion
used in \cite[lines~1278--1282]{GrossNesmeVogtsWerner2012}. It must not be
cited as (\ref{eq:sw04_csform_matrix_scope_source_decomposition}): neither the
abstract $C^\ast$-algebra generality nor the decomposition
(\ref{eq:sw04_csform_matrix_scope_central_cut}) into minimal central
projections is formalized. Removing the restriction requires the minimal
central projections of $S$ and a compatible reassembly of the summands; the
building block proved here is the step that reassembly would repeat for each
$\mu$.\footnote{Consumer: \ghissue{463}.}

\bibliographystyle{alpha}
\bibliography{references}

\end{document}
